\section{Related Work}
\textbf{Vision Pretraining.}
Pretraining ConvNets~\cite{krizhevsky2012imagenet} or Transformers~\cite{vaswani2017attention} on large-scale annotated data such as ImageNet~\cite{girshick2014rich, long2015fully, simonyan2014two}, Instagram~\cite{mahajan2018exploring} or JFT~\cite{zhai2021scaling} has become a popular strategy towards solving visual recognition problems including classification, localization, segmentation, video recognition, tracking and many other problems. Recently, self-supervised pretraining approaches have also been explored. BEiT~\cite{bao2021beit} proposes a masked image modeling task following BERT~\cite{devlin2018bert} in natural language processing, and uses quantized visual token ids as prediction targets. MAE~\cite{he2021masked} and SimMIM~\cite{xie2021simmim} remove the need for an image tokenizer and directly use a light-weight decoder or projection layer to regress pixel values.
Nonetheless,
these methods only learn models for the vision modality and thus they are not applicable to tasks that require joint reasoning over both image and text inputs.

\textbf{Vision-Language Pretraining.}
In recent years, rapid progress has been made in vision-language pretraining (VLP),
which aims to jointly encode vision and language in a fusion model.
Early work (e.g. LXMERT \cite{tan2019lxmert},
UNITER \cite{chen2020uniter}, VinVL \cite{Zhang_2021_CVPR}) in this direction relies on pretrained object detection modules such as Fast(er) R-CNN \citep{NIPS2015_14bfa6bb} to extract visual representations.
Later efforts such as ViLT \cite{kim2021vilt} and VLMo \cite{wang2021vlmo} unify vision and language transformers,
and train a multimodal transformer from scratch.

\textbf{Image-Text Foundation Models.}
Recent work has proposed image-text foundation models that can subsume both vision and vision-language pretraining.
CLIP~\cite{radford2021learning} and ALIGN~\cite{jia2021scaling} demonstrate that dual-encoder models pretrained with contrastive objectives on noisy image-text pairs 
can learn strong image and text representations for crossmodal alignment tasks and zero-shot image classification.
Florence~\cite{yuan2021florence} further develops this method with unified contrastive objective~\cite{yang2022unified}, training foundation models that can be adapted for a wide range of vision and image-text benchmarks. To further improve zero-shot image classification accuracy, LiT~\cite{zhai2021lit} and BASIC~\cite{pham2021combined} first pretrain model on an large-scale image annotation dataset with cross-entropy and further finetune with contrastive loss on an noisy alt-text image dataset.
Another line of research~\cite{wang2021simvlm, wang2022unifying, piergiovanni2022answer} proposes encoder-decoder models trained with generative losses and shows strong results in vision-language benchmarks while the visual encoder still performs competitively on image classification. 
In this work, we focus on training an image-text foundation model from scratch in a single pretraining stage to unify these approaches.
While recent works ~\cite{singh2021flava,li2021align,li2022blip} have also explored image-text unification,
they require multiple pretraining stages of unimodal and multimodal modules to attain good performance.
For example, ALBEF~\cite{li2021align} combines contrastive loss with masked language modelling (MLM) with a dual-encoder design. 
However, our approach is simpler and more efficient to train while also enables more model capabilities:
(1) CoCa only performs one forward and backward propagation for a batch of image-text pairs while ALBEF requires two (one on corrupted inputs and another without corruption),
(2) CoCa is trained from scratch on the two objectives only while ALBEF is initialized from pretrained visual and textual encoders with additional training signals including momentum modules.
(3) The decoder architecture with generative loss is preferred for natural language generation and thus directly enables image captioning and zero-shot learning~\cite{wang2021simvlm}.